%% 02b-lalonde.tex — Glasgow LaLonde walkthrough: history + four estimators
%% Scott Cunningham, University of Glasgow, May 2026
%% Adapted from the Zurich Day 1 LaLonde deck.

\documentclass[11pt,aspectratio=169]{beamer}
\usepackage{remix}
\usepackage{booktabs}
\usepackage{listings}
\usepackage{colortbl}
\renewcommand{\remixfooter}{University of Glasgow\,\textperiodcentered\,May 2026}

\begin{document}

%% ---------- Title ----------
\remixtitleframe
  {LaLonde, Dehejia--Wahba, and Why the 2$\times$2 Still Matters}
  {Selection on observables, in the case that built the literature}
  {Scott Cunningham \textperiodcentered\ Baylor University}
  {University of Glasgow \textperiodcentered\ May 2026}

%% =====================================================================
%%                                LESSON 1
%% =====================================================================
\remixsection{1.\ LaLonde 1986}

\begin{frame}{The National Supported Work Demonstration}
\vspace{1em}
\begin{center}
\large
A federally funded \textbf{randomized} job-training program.
\end{center}
\vspace{0.4em}
\begin{center}
\large
1975--1979. Four target groups, including \\[0.2em] disadvantaged adult men.
\end{center}
\vspace{0.6em}
\begin{center}
\small\color{warmgray}\itshape
Random assignment to treatment vs.\ control. \;Outcome: real earnings in 1978.
\end{center}
\end{frame}

\begin{frame}{LaLonde's experimental benchmark}
\vspace{1.6em}
\begin{center}
\Large
$\widehat{\text{ATT}}^{\text{RCT}} \;\approx\; \$886$
\end{center}
\vspace{0.6em}
\begin{center}
1978 dollars. \;Adult-male subsample.
\end{center}
\vspace{1em}
\begin{center}
\small\color{warmgray}\itshape
Random assignment $\Rightarrow$ no selection. \;This is the \textit{truth}.
\end{center}
\end{frame}

\begin{frame}{Then LaLonde threw away the controls}
\vspace{1em}
\begin{center}
\large
He kept the treated --- the 297 NSW men.
\end{center}
\vspace{0.6em}
\begin{center}
\large
He replaced the controls with samples from the \\[0.2em] \textbf{CPS} and the \textbf{PSID}.
\end{center}
\vspace{0.8em}
\begin{center}
\small\color{warmgray}\itshape
A simulation of what an applied researcher would face --- \\
no experiment, just observational data and methods.
\end{center}
\end{frame}

\begin{frame}{Then he tried every method going at the time}
\vspace{0.6em}
\begin{itemize}
  \item OLS with controls
  \item Two-step Heckman selection corrections
  \item Fixed effects on pre-program earnings
  \item Difference-in-differences
\end{itemize}
\vspace{0.6em}
\begin{center}
\large\itshape\color{charcoal}
Could any of them recover \$886?
\end{center}
\end{frame}

\begin{frame}{No}
\vspace{1.6em}
\begin{center}
\large
Estimates ranged from \;\highlight{$-$\$16{,}000} \;to \;\highlight{$+$\$3{,}000}.
\end{center}
\vspace{1em}
\begin{center}
\large
Most were negative. \;Most missed the benchmark by orders of magnitude.
\end{center}
\vspace{1em}
\begin{center}
\small\color{warmgray}\itshape
LaLonde (1986), \textit{AER} 76(4): 604--620.
\end{center}
\end{frame}

\begin{frame}{The headline that shook applied econometrics}
\vspace{1.6em}
\begin{center}
\Large\itshape\color{charcoal}
``The econometric procedures \\[0.2em] cannot replicate the experimental results.''
\end{center}
\vspace{1.4em}
\begin{center}
\small\color{warmgray}\itshape
This is the paper that made program evaluation \\
move toward randomization and design.
\end{center}
\end{frame}

%% =====================================================================
%%                                LESSON 2
%% =====================================================================
\remixsection{2.\ Dehejia \& Wahba 2002}

\begin{frame}{Dehejia \& Wahba reopened the case}
\vspace{1em}
\begin{center}
\large
\textbf{Rajeev Dehejia} and \textbf{Sadek Wahba}, \\[0.2em] \textit{Review of Economics and Statistics} (2002).
\end{center}
\vspace{0.8em}
\begin{center}
\large
They asked: \;does \textbf{propensity score matching} \\[0.2em] do better than LaLonde's methods?
\end{center}
\end{frame}

\begin{frame}{First, they re-cut the experimental sample}
\vspace{1em}
\begin{center}
\large
They restricted to NSW men with earnings observed in \\[0.2em] \textbf{both 1974 and 1975}.
\end{center}
\vspace{0.6em}
\begin{center}
\large
A smaller, slightly different population.
\end{center}
\vspace{1em}
\begin{center}
\small\color{warmgray}\itshape
This is the subsample the homework uses: \texttt{lalonde\_exp\_panel.dta}.
\end{center}
\end{frame}

\begin{frame}{A new experimental benchmark}
\vspace{1.6em}
\begin{center}
\Large
$\widehat{\text{ATT}}^{\text{RCT}}_{\text{D\&W}} \;\approx\; \$1{,}794$
\end{center}
\vspace{0.6em}
\begin{center}
\large
Roughly \;\highlight{\$1{,}700} \;in the way we usually report it.
\end{center}
\vspace{1em}
\begin{center}
\small\color{warmgray}\itshape
Why bigger than LaLonde's \$886? \;\textit{Different subsample, different ATT.}
\end{center}
\end{frame}

\begin{frame}{Why the headlines differ}
\vspace{0.5em}
\begin{itemize}
  \item Both numbers come from the \textit{same} NSW experiment.
  \item LaLonde used the full adult-male sample.
  \item D\&W used a \textit{subset} --- only men with two years of pre-program earnings.
  \item Subsetting on a covariate changes the population whose ATT you're estimating.
  \item Same identification (random assignment). \;Different target group.
\end{itemize}
\vspace{0.5em}
\begin{center}
\small\color{warmgray}\itshape
This is not a contradiction --- it is the ATT defined on different units.
\end{center}
\end{frame}

\begin{frame}{D\&W's quasi-experimental result}
\vspace{1em}
\begin{center}
\large
With \textbf{propensity-score matching} on \\[0.2em] CPS / PSID controls...
\end{center}
\vspace{0.6em}
\begin{center}
\Large
they recovered \;\highlight{$\approx \$1{,}700$}.
\end{center}
\vspace{0.8em}
\begin{center}
\small\color{warmgray}\itshape
A constructive answer to LaLonde: \;design plus matching \textit{can} work.
\end{center}
\end{frame}

%% =====================================================================
%%                                LESSON 3
%% =====================================================================
\remixsection{3.\ Why the ATT didn't move}

\begin{frame}{A puzzle}
\vspace{1em}
\begin{center}
\large
We used \textbf{non-experimental controls} (CPS, PSID) \\[0.2em] in the larger LaLonde--D\&W literature.
\end{center}
\vspace{0.6em}
\begin{center}
\large
And yet the ATT we target is still \;\highlight{$\approx \$1{,}700$.}
\end{center}
\vspace{0.8em}
\begin{center}
\large\itshape\color{charcoal}
Why doesn't the ATT change when we change the controls?
\end{center}
\end{frame}

\begin{frame}{Because the ATT is a property of the treated}
\vspace{1.4em}
\begin{center}
\Large
$\text{ATT} \;=\; \mathbb{E}[Y^1 - Y^0 \mid D = 1,\, \text{Post}]$
\end{center}
\vspace{1.2em}
\begin{center}
\large
Only the \;\highlight{$D = 1$} \;side appears.
\end{center}
\vspace{0.6em}
\begin{center}
The control units are just a \textit{tool} to learn about $\mathbb{E}[Y^0 \mid D{=}1, \text{Post}]$.
\end{center}
\end{frame}

\begin{frame}{Random assignment gave us a clean tool}
\vspace{1.4em}
\begin{center}
\large
Under random assignment:
\end{center}
\vspace{0.4em}
\begin{center}
\Large
$\mathbb{E}[Y^0 \mid D{=}1] \;=\; \mathbb{E}[Y^0 \mid D{=}0]$
\end{center}
\vspace{1em}
\begin{center}
So $\mathbb{E}[Y \mid D{=}0, \text{Post}]$ \textit{is} the counterfactual we need. \\[0.2em]
No selection bias.
\end{center}
\end{frame}

\begin{frame}{Non-experimental controls break that equality}
\vspace{1em}
\begin{center}
\large
CPS men and NSW participants are not interchangeable.
\end{center}
\vspace{0.6em}
\begin{center}
\large
$\mathbb{E}[Y^0 \mid D{=}1] \;\neq\; \mathbb{E}[Y^0 \mid D{=}0]$
\end{center}
\vspace{1em}
\begin{center}
\large
Selection bias enters.
\end{center}
\vspace{0.6em}
\begin{center}
\small\color{warmgray}\itshape
NSW men had worse pre-program earnings. \;The CPS pool didn't.
\end{center}
\end{frame}

\begin{frame}{But the \textit{target} hasn't changed}
\vspace{1.2em}
\begin{center}
\large
$\mathbb{E}[Y^0 \mid D{=}1, \text{Post}]$ is a feature of the NSW men.
\end{center}
\vspace{0.6em}
\begin{center}
\large
It does not depend on who we picked as controls.
\end{center}
\vspace{1em}
\begin{center}
\large
What the controls control is \textit{our estimate} of that target --- \\[0.2em] and the size of the bias.
\end{center}
\end{frame}

\begin{frame}{So what is DiD for, in this story?}
\vspace{0.6em}
\begin{itemize}
  \item In the RCT sample: DiD just \textit{confirms} the post-period difference-in-means.
  \item In the non-experimental sample: DiD tries to \textit{net out} the pre-existing level gap between NSW and CPS men.
  \item The truth is fixed at \$1{,}700. \;The question is whether DiD's parallel-trends assumption is enough to recover it.
\end{itemize}
\vspace{0.6em}
\begin{center}
\large\itshape\color{charcoal}
Did it?
\end{center}
\end{frame}

\begin{frame}{The benchmark and the challenge}
\vspace{1em}
\begin{center}
\Large
RCT benchmark: \;\highlight{$\$1{,}794$}
\end{center}
\vspace{0.6em}
\begin{center}
\Large
Naive DiD on the non-experimental panel: \;\highlight{$\$3{,}621$}
\end{center}
\vspace{0.8em}
\begin{center}
\large
We can do better than additive controls. \;Let's go to the code.
\end{center}
\end{frame}

%% =====================================================================
%%                                LESSON 5
%% =====================================================================
\remixsection{4.\ Now make it harder}

\begin{frame}{Throw away the experimental controls}
\vspace{1em}
\begin{center}
\large
Keep the 297 NSW treated men.
\end{center}
\vspace{0.6em}
\begin{center}
\large
Replace the controls with a random sample of \\[0.2em] American households from the \textbf{CPS}.
\end{center}
\vspace{1em}
\begin{center}
\small\color{warmgray}\itshape
Data: \texttt{lalonde\_nonexp\_panel.dta}. \;Same treated rows. \;New controls.
\end{center}
\end{frame}

\begin{frame}{Selection is severe}
\vspace{0.8em}
\begin{center}
\small
\renewcommand{\arraystretch}{1.3}
\begin{tabular}{@{}lrrr@{}}
\toprule
        & \textbf{1974}    & \textbf{1975}    & \textbf{1978}    \\
\midrule
NSW (treated, $n=185$)   & \$2{,}096        & \$1{,}532        & \$6{,}349        \\
CPS (control, $n=15{,}992$) & \$14{,}017       & \$13{,}651       & \$14{,}847       \\
\midrule
Diff                      & $-\$11{,}921$    & $-\$12{,}119$    & $-\$8{,}498$    \\
\bottomrule
\end{tabular}
\end{center}
\vspace{0.5em}
\begin{center}
\large
$\mathbb{E}[Y^0 \mid D{=}1] \;\ll\; \mathbb{E}[Y^0 \mid D{=}0]$
\end{center}
\vspace{0.3em}
\begin{center}
\small\color{warmgray}\itshape
The level gap is huge --- and so is the \textit{growth-rate} gap. \;DiD doesn't fix that.
\end{center}
\end{frame}

%% ---------- Animated 2x2: show DiD arithmetic on the same table ----------
\begin{frame}{The 2$\times$2: where \$3{,}621 comes from}
\vspace{0.5em}
\begin{center}
\small
\renewcommand{\arraystretch}{1.5}
\begin{tabular}{@{}lrrr@{}}
\toprule
        & \textbf{1974}    & \textbf{1975 (pre)}    & \textbf{1978 (post)}    \\
\midrule
NSW (treated, $n=185$)   & \$2{,}096   & \cellcolor{accent!25}\$1{,}532    & \cellcolor{accent!25}\$6{,}349 \\
CPS (control, $n=15{,}992$) & \$14{,}017 & \cellcolor{ocean!25}\$13{,}651 & \cellcolor{ocean!25}\$14{,}847 \\
\bottomrule
\end{tabular}
\end{center}

\vspace{0.6em}
\begin{center}
\small
\begin{tabular}{@{}r@{\,\,}c@{\,\,}r@{\,\,}c@{\,\,}r@{}}
\textcolor{accent}{\textbf{NSW change:}}   & & \textcolor{accent}{\$6{,}349 $-$ \$1{,}532}       & $=$ & \textcolor{accent}{\$4{,}817} \\
\textcolor{ocean}{\textbf{CPS change:}}    & & \textcolor{ocean}{\$14{,}847 $-$ \$13{,}651}      & $=$ & \textcolor{ocean}{\$1{,}196} \\
\midrule
\textbf{DiD}                                & & \textcolor{accent}{\$4{,}817} $-$ \textcolor{ocean}{\$1{,}196} & $=$ & $\boxed{\$3{,}621}$ \\
\end{tabular}
\end{center}

\vspace{0.6em}
\begin{center}
\Large
\highlight{$\$3{,}621 \;\neq\; \$1{,}794$}
\end{center}
\vspace{0.2em}
\begin{center}
\small\color{warmgray}\itshape
NSW men bounce back fast from a deep dip. \;CPS men don't. \;DiD attributes the entire gap in growth rates to the program.
\end{center}
\end{frame}

\begin{frame}{Naive DiD on this sample}
\vspace{1.2em}
\begin{center}
\Large
$\widehat{\delta}_{\text{naive DiD}} \;=\; +\$3{,}621$
\end{center}
\vspace{0.6em}
\begin{center}
\large
\textit{Right sign. \;Wrong magnitude.} \;The benchmark is \;\highlight{$\$1{,}794$}.
\end{center}
\vspace{0.8em}
\begin{center}
\small\color{warmgray}\itshape
NSW men had unusually low pre-earnings ($\$1{,}532$) and bounced back fast. \\
CPS controls didn't. \;DiD attributes the gap in growth rates entirely to treatment.
\end{center}
\end{frame}

%% =====================================================================
%%                                LESSON 6
%% =====================================================================
\remixsection{5.\ Run the specs}

%% ---------- Spec 0 ----------
\begin{frame}[fragile]{Spec 0: Unconditional parallel trends, no covariates}
\textbf{Assumption.} \;Naive parallel trends: $E[\Delta Y(0) \mid D{=}1] = E[\Delta Y(0) \mid D{=}0]$ unconditionally.
\vspace{0.6em}

\textbf{Run:}
\begin{lstlisting}[language={[ANSI]C},basicstyle=\ttfamily\small]
reg re i.post##i.ever_treated, robust
\end{lstlisting}

\vspace{0.4em}
\textbf{Compare to \$1,794.}
\end{frame}

%% ---------- Spec A ----------
\begin{frame}[fragile]{Spec A: Conditional PT, additive $X$}
\textbf{Assumption.} \;Conditional parallel trends with one slope per covariate (pre = post).
\vspace{0.6em}

\textbf{Run:}
\begin{lstlisting}[language={[ANSI]C},basicstyle=\ttfamily\small]
global X "age agesq educ educsq marr nodegree black hisp re74 re75 u74 u75"

reg re i.post##i.ever_treated $X, robust
\end{lstlisting}

\vspace{0.4em}
\textbf{Compare to \$1,794.}
\end{frame}

%% ---------- Spec B ----------
\begin{frame}[fragile]{Spec B: Conditional PT, $\text{Post}\times X$}
\textbf{Assumption.} \;Conditional parallel trends, allowing $X$-effects to differ pre vs post. \;Use \texttt{\#\#} so the spec includes \textit{both} the pre-period $X$ effect and its post-period change.
\vspace{0.4em}

\textbf{Run:}
\begin{lstlisting}[language={[ANSI]C},basicstyle=\ttfamily\footnotesize]
reg re i.post##i.ever_treated ///
    i.post##c.age i.post##c.agesq i.post##c.educ i.post##c.educsq ///
    i.post##i.marr i.post##i.nodegree i.post##i.black i.post##i.hisp ///
    i.post##c.re74 i.post##c.re75 i.post##i.u74 i.post##i.u75, robust
\end{lstlisting}

\vspace{0.3em}
\textbf{Compare to \$1,794.} \;\meta{Single-\texttt{\#} forces pre-period $X$ effect to zero --- not what we want.}
\end{frame}

%% ---------- Spec C ----------
\begin{frame}[fragile]{Spec C: Saturated TWFE / OR (HIT 1997)}
\textbf{Assumption.} \;Conditional parallel trends; impute $Y(0)$'s trend on controls only and average over treated.
\vspace{0.4em}

\textbf{Run --- by hand:}
\begin{lstlisting}[language={[ANSI]C},basicstyle=\ttfamily\footnotesize]
reshape wide re, i(id) j(post)
gen dy = re1 - re0

reg dy $X if ever_treated == 0
predict dy_hat
gen delta_hat = dy - dy_hat if ever_treated == 1
summarize delta_hat if ever_treated == 1
\end{lstlisting}

\textbf{Cross-check with \texttt{drdid}:}
\begin{lstlisting}[language={[ANSI]C},basicstyle=\ttfamily\footnotesize]
drdid re $X, time(year) ivar(id) tr(ever_treated) reg
\end{lstlisting}

\textbf{Compare to \$1,794.} \;\meta{Two numbers should agree to the dollar.}
\end{frame}

%% ---------- Spec D ----------
\begin{frame}[fragile]{Spec D: IPW (Abadie 2005)}
\textbf{Assumption.} \;Conditional PT; correct propensity score. \;Treated weighted by 1; controls weighted by $\hat p(X)/(1-\hat p(X))$.
\vspace{0.3em}

\textbf{Run --- by hand:}
\begin{lstlisting}[language={[ANSI]C},basicstyle=\ttfamily\scriptsize]
logit ever_treated $X
predict phat, pr
quietly sum ever_treated
local pr_treat = r(mean)
gen w_ipw = (ever_treated - phat) / (1 - phat) / `pr_treat'
gen ipw_dy = w_ipw * dy
summarize ipw_dy
\end{lstlisting}

\textbf{Cross-check with \texttt{drdid}:}
\begin{lstlisting}[language={[ANSI]C},basicstyle=\ttfamily\footnotesize]
drdid re $X, time(year) ivar(id) tr(ever_treated) ipw
\end{lstlisting}

\textbf{Compare to \$1,794.}
\end{frame}

%% ---------- Spec E ----------
\begin{frame}[fragile]{Spec E: DR-DiD (Sant'Anna \& Zhao 2020)}
\textbf{Assumption.} \;Conditional PT, AND \textit{either} $\hat\mu_0$ or $\hat p(X)$ is correct (not both required).
\vspace{0.3em}

\textbf{Run --- by hand (in wide-FD form, \texttt{phat}/\texttt{dy\_hat} already computed):}
\begin{lstlisting}[language={[ANSI]C},basicstyle=\ttfamily\scriptsize]
gen dr_t = ever_treated * (dy - dy_hat) / `pr_treat'
gen dr_c = (1 - ever_treated) * (phat/(1-phat)) * (dy - dy_hat) / `pr_treat'
quietly sum dr_t
local mt = r(mean)
quietly sum dr_c
display "DR ATT = " %9.0fc `mt' - r(mean)
\end{lstlisting}

\textbf{Cross-check with \texttt{drdid}:}
\begin{lstlisting}[language={[ANSI]C},basicstyle=\ttfamily\footnotesize]
drdid re $X, time(year) ivar(id) tr(ever_treated) dripw
\end{lstlisting}

\textbf{Compare to \$1,794.}
\end{frame}

%% =====================================================================
%%                                LESSON 7
%% =====================================================================
\remixsection{6.\ How close did each one get?}

\begin{frame}{The comparison}
\vspace{0.2em}
\begin{center}
\scriptsize
\renewcommand{\arraystretch}{1.15}
\begin{tabular}{@{}lrr@{}}
\toprule
\textbf{Spec} & \textbf{ATT} & \textbf{Distance from \$1{,}794} \\
\midrule
0. Naive TWFE (no $X$)                              & $+\$3{,}621$ & $+\$1{,}827$ \\
A. Additive $X$                                     & $+\$3{,}621$ & $+\$1{,}827$ \\
B. Post${\times}X$ (\texttt{\#\#} syntax)            & $+\$1{,}184$ & $-\$610$ \\
C. Saturated TWFE (FD)                              & $+\$1{,}263$ & $-\$531$ \\
\;\;OR by hand $=$ \texttt{drdid reg}                & $+\$1{,}263$ & $-\$531$ \\
D. IPW by hand $=$ \texttt{drdid ipw}                & $+\$1{,}361$ & $-\$433$ \\
E. DR by hand                                       & $+\$1{,}496$ & $-\$298$ \\
\;\;DR \texttt{drdid dripw}                          & $+\$1{,}485$ & $-\$309$ \\
\midrule
\textbf{RCT benchmark}                              & \textbf{$+\$1{,}794$} & --- \\
\bottomrule
\end{tabular}
\end{center}
\vspace{0.3em}
\begin{center}
\footnotesize\color{warmgray}\itshape
Naive $=$ Additive-$X$: the linear $X$ control is absorbed. \;DR closest at \$298 below. \;By-hand $=$ \texttt{drdid} to the dollar for OR and IPW; DR within \$11.
\end{center}
\end{frame}

\begin{frame}{Seven specs, one benchmark}
\begin{center}
\includegraphics[height=0.70\textheight]{figures/lalonde_forest.png}
\end{center}
\begin{center}
\scriptsize\color{warmgray}\itshape
Naive and Additive-$X$ CIs sit entirely above \$1{,}794. \;Every design-aware spec contains it.
\end{center}
\end{frame}

\begin{frame}{DR lands closest; the others bracket it}
\vspace{0.5em}
\begin{itemize}
  \item \textbf{DR} (\$1{,}496) sits \textbf{\$298 below} --- closest observational estimator.
  \item \textbf{IPW} (\$1{,}361) sits \$433 below.
  \item \textbf{OR / Saturated TWFE} (\$1{,}263) sits \$531 below --- two paths, same number.
  \item \textbf{Post${\times}X$} (\$1{,}184) sits \$610 below --- still misses Bias$_\text{HTE}$.
  \item \textbf{Additive $X$} (\$3{,}621) doesn't move off Naive (\$3{,}621) --- the FE absorbed the linear $X$.
\end{itemize}
\vspace{0.4em}
\begin{center}
\small\color{warmgray}\itshape
``Doubly robust'' = robust to one misspecification, not best in every sample. \;DR's $298$ below is selection on unobservables --- the residual no observational $X$ can close.
\end{center}
\end{frame}

\begin{frame}{What this lab teaches you}
\vspace{0.8em}
\begin{itemize}
  \item The \textbf{ATT is fixed} ($\approx +\$1{,}794$). \;The controls only change the bias.
  \item \textbf{Additive $X$ barely helps:} \$3{,}621 vs naive \$3{,}621. \;You need a design-aware estimator.
  \item Post${\times}X$, OR, IPW, DR offer four different uses of the same $X$, all closer to truth.
  \item On this dataset, \textbf{DR-DiD lands closest} at \$298 below \$1{,}794.
  \item No method is magic. Each rests on an assumption you can name.
\end{itemize}
\end{frame}

%% =====================================================================
%%   Closing synthesis: it was a TWFE specification problem all along
%% =====================================================================

\begin{frame}{What additive $X$ actually assumed}
\vspace{0.4em}
The additive-$X$ spec\, $Y_{it} = \alpha + \delta(D_i{\cdot}\text{post}_t) + X_i'\beta + \varepsilon$\, was a \textbf{TWFE specification choice}. \;Two restrictions came along free of charge:
\vspace{0.6em}
\begin{enumerate}\setlength\itemsep{4pt}
  \item \textbf{Homogeneous treatment effect in $X$.} \;One scalar $\delta$. \;No room for $\delta(X)$.
  \item \textbf{No time-varying $X$ effect.} \;One $\beta$, used in both pre and post. \;If $X$'s relationship with $Y(0)$ shifts between periods, the spec can't see it.
\end{enumerate}
\vspace{0.6em}
\highlight{Neither restriction is identification. \;Both are specification.}
\end{frame}

\begin{frame}{Two TWFE relaxations}
\vspace{0.4em}
We saw both in the synthetic simulation. \;Both are still TWFE --- just less restrictive parameterizations.
\vspace{0.6em}

\textbf{Post $\times X$.} \;Let $\beta$ differ pre vs post:
\[
Y = \alpha_1 + \alpha_2 T + \alpha_3 D + \delta(DT) + \beta_1 X + \beta_2 XT + \varepsilon
\]
Relaxes restriction (2). \;Still imposes (1): one scalar $\delta$.

\vspace{0.5em}
\textbf{Fully saturated.} \;Let $\beta$ differ pre/post \textit{and} let $\delta$ depend on $X$:
\[
Y = \alpha_1 + \alpha_2 T + \alpha_3 D + \delta(DT) + \beta_1 X + \beta_2 XT + \beta_3 XD + \beta_4 DTX + \varepsilon
\]
Relaxes both. \;The ATT is now $\delta + \beta_4\,\mathbb{E}[X \mid D{=}1]$ --- average the $X$-specific effect over the treated.
\end{frame}

\begin{frame}{The punchline}
\vspace{0.8em}
\begin{center}
\large
Additive $X$ failed on LaLonde \\[0.3em]
not because the data are weird, \\[0.3em]
but because the \highlight{TWFE specification} was too restrictive.
\end{center}

\vspace{0.8em}
\begin{center}
\large
\textbf{Fully saturated TWFE} = the regression form of HIT 1997 outcome regression. \\[0.3em]
Same estimator. \;Same ATT. \;Different machinery.
\end{center}

\vspace{0.6em}
\begin{center}
\small\color{warmgray}\itshape
Spec C (FD with $D \times X$ interactions) and OR-by-hand returned \textit{the same dollar:} \$1{,}263.
\end{center}
\end{frame}

%% =====================================================================
%%                                CLOSING
%% =====================================================================
\begin{frame}[plain]
\begin{tikzpicture}[remember picture,overlay]
\fill[cream] (current page.south west) rectangle (current page.north east);
\node[anchor=center, text width=0.84\paperwidth, align=center, text=charcoal] at (current page.center) {%
  \fontsize{20pt}{30pt}\selectfont
  The ATT lives on the \textbf{treated.}\\[0.6em]
  The controls only tell you \\[0.2em]
  whether your method \\[0.2em]
  can {\color{accent}\bfseries see it.}
};
\draw[accent, line width=1pt] ([yshift=1.8cm,xshift=0.28\paperwidth]current page.south west) -- ([yshift=1.8cm,xshift=-0.28\paperwidth]current page.south east);
\end{tikzpicture}
\end{frame}

\end{document}
